\chapter{Lifted spectral covers, tensor defect, and cubic resolvent towers}
\label{ch:lifted-spectral-defect}

The visible generating function is only a shadow.  What controls transport,
branching, and tensor composition is a finite cover over the spectral line
$X=\mathbf P^1_x$ together with the rational mechanism by which the visible
series descends from that cover.  In the quadratic examples of
Theorem~\ref{thm:ds-bar-gf-discriminant} the tautological cover generated by
$P_{\cA^!}(x)$ itself already sees the relevant geometry.  In rational-shadow
regimes such as the conjectural
$\widehat{\mathfrak{sl}}_3/\mathcal W_3$ orbit
(Conjectures~\ref{conj:sl3-bar-gf} and~\ref{conj:w3-bar-gf}), the tautological
field $k(x)(P_{\cA^!}(x))$ collapses to $k(x)$ and therefore forgets the
nontrivial transport data.  The correct object is a \emph{lifted spectral
presentation}.  This chapter constructs that object from first principles and
builds an exact defect calculus in the quadratic and cubic $S_3$ regimes.

Throughout this chapter $k$ is an algebraically closed field of
characteristic $0$; when cubic Kummer theory appears we also assume
$\mu_3\subset k$.  Set
\[
F:=k(x), \qquad X:=\mathbf P^1_x.
\]
All covers are normalizations of $X$ in finite separable extensions of~$F$.

\section{Spectral presentations and their canonical branch packages}
\label{sec:spectral-presentations}

\begin{definition}[Spectral presentation]
\label{def:spectral-presentation}
Let $P(x)\in \overline{F}$ be a visible spectral series.
A \emph{spectral presentation} of~$P$ consists of
\begin{enumerate}[label=\textup{(\roman*)}]
\item a finite separable extension $L/F$,
\item a primitive element $y\in L$ with $L=F(y)$,
\item a rational reconstruction law $r(x,t)\in F(t)$ such that
\[
P(x)=r\bigl(x,y\bigr)\in L.
\]
\end{enumerate}
The presentation is \emph{tautological} if $y=P$ and $r(x,t)=t$.
It is \emph{hidden} if $P\in F$ but $L\neq F$.
\end{definition}

\begin{remark}[Why the presentation is the right object]
\label{rem:presentation-right-object}
Theorem~\ref{thm:discriminant-spectral} attaches to a visible series a
spectral shadow: a denominator in the rational case, or a discriminant in the
algebraic case.  That shadow records branch values but not the full sheet.  In
particular, if $P\in F$ is rational then the tautological field is trivial, even
though the same visible series may arise as a rational expression in a
nontrivial quadratic or cubic coordinate on a finite cover.  Theorems
\ref{thm:quadratic-common-sheet-law}--\ref{thm:cubic-rational-shadow-frontier}
show that tensor transport and DS transport depend on the lifted cover, not
merely on the visible shadow.
\end{remark}

\begin{definition}[Equivalence of presentations]
\label{def:equiv-presentations}
Two spectral presentations $(L,y,r)$ and $(L',y',r')$ of the same visible
series $P$ are \emph{birationally equivalent} if there exist rational functions
$u(t),v(t)\in F(t)$ such that
\[
y'=u(y),\qquad y=v(y').
\]
Equivalently, $F(y)=F(y')$ as subfields of $\overline{F}$.
\end{definition}

\begin{proposition}[Birational invariance of the cover; \ClaimStatusProvedHere]
\label{prop:birational-invariance-cover}
Birationally equivalent presentations determine the same normalized spectral
cover over~$X$.
\end{proposition}

\begin{proof}
If $y'=u(y)$ and $y=v(y')$ with $u,v\in F(t)$, then
$F(y)=F(y')=:L$.  Normalization of~$X$ depends only on the function field, so
both presentations determine the normalization of $X$ in~$L$.
\end{proof}

\begin{construction}[Lifted spectral cover package]
\label{cons:lifted-package}
Let $(L,y,r)$ be a spectral presentation of~$P$ and let
\[
\pi\colon \Sigma\to X
\]
be the normalization of $X$ in~$L$.  The \emph{lifted spectral cover package}
consists of the following objects.
\begin{enumerate}[label=\textup{(\roman*)}]
\item The \emph{sheet field} $L=F(y)$ and the associated
$F$-vector space $E:=L$.
\item The \emph{sheet operator}
$M_y\in \operatorname{End}_F(E)$ given by multiplication by~$y$.
\item The \emph{branch complex}
\[
\mathbb B(\Sigma/X):=L_{\Sigma/X}.
\]
\item The \emph{branch space}
\[
V^{\mathrm{br}}(\Sigma/X):=H^0\bigl(\Sigma,\Omega^1_{\Sigma/X}\bigr).
\]
\item If $\pi$ is unramified over $x=0$, the \emph{branch operator}
\[
T^{\mathrm{br}}_{\Sigma/X}:=\pi^*(1/x)\cdot \mathrm{id}
\quad\text{on }V^{\mathrm{br}}(\Sigma/X).
\]
\end{enumerate}
\end{construction}

\begin{proposition}[Minimal polynomial and sheet rank; \ClaimStatusProvedHere]
\label{prop:sheet-operator-minpoly}
Let $(L,y,r)$ be a spectral presentation.  Let $m_y(T)\in F[T]$ be the monic
minimal polynomial of~$y$ over~$F$.  Then:
\begin{enumerate}[label=\textup{(\roman*)}]
\item the $F[T]$-module $E$ is cyclic with
$E\cong F[T]/(m_y)$;
\item the characteristic polynomial of $M_y$ is $m_y(T)$;
\item the sheet rank is
\[
\dim_F E=[L:F]=\deg_T m_y.
\]
\end{enumerate}
\end{proposition}

\begin{proof}
Because $L=F(y)$, the map $F[T]\to L$ sending $T\mapsto y$ is surjective with
kernel $(m_y)$, proving~(i).  In the basis
$1,y,\dots,y^{d-1}$ with $d=\deg m_y$, the operator $M_y$ is represented by the
companion matrix of $m_y$, so its characteristic polynomial is exactly~$m_y$.
Statement~(iii) follows immediately.
\end{proof}

\begin{proposition}[Support of the branch complex; \ClaimStatusProvedHere]
\label{prop:branch-complex-support}
The complex $\mathbb B(\Sigma/X)$ is perfect of amplitude $[-1,0]$, and
$H^0\bigl(\mathbb B(\Sigma/X)\bigr)=\Omega^1_{\Sigma/X}$ is a torsion coherent
sheaf supported on the ramification locus of~$\pi$.  In particular,
$V^{\mathrm{br}}(\Sigma/X)$ is finite-dimensional over~$k$.
\end{proposition}

\begin{proof}
The morphism $\pi$ is finite between regular curves in characteristic~$0$.
Hence $L_{\Sigma/X}$ has Tor-amplitude $[-1,0]$.  Since finite separable maps of
curves are generically étale, the sheaf $\Omega^1_{\Sigma/X}$ vanishes away from
the ramification locus and is therefore torsion of finite length.  Properness of
$\Sigma$ over~$k$ makes $H^0$ finite-dimensional.
\end{proof}

\begin{theorem}[Determinantal branch formula; \ClaimStatusProvedHere]
\label{thm:determinantal-branch-formula-general}
Assume $\pi\colon \Sigma\to X$ is unramified over $x=0$.  Let
$\operatorname{Supp}\Omega^1_{\Sigma/X}=\{p_1,\dots,p_N\}$ and set
\[
\ell_i:=\operatorname{length}_{\mathcal O_{\Sigma,p_i}}
\Omega^1_{\Sigma/X,p_i}.
\]
Then
\[
\det\bigl(1-zT^{\mathrm{br}}_{\Sigma/X}\bigr)
=
\prod_{i=1}^{N}\left(1-\frac{z}{x(\pi(p_i))}\right)^{\ell_i}.
\]
In particular, the branch operator records the branch divisor with its full
scheme-theoretic multiplicities.
\end{theorem}

\begin{proof}
Since $\Omega^1_{\Sigma/X}$ is a finite-length sheaf on a proper curve,
\[
V^{\mathrm{br}}(\Sigma/X)
\cong
\bigoplus_{i=1}^{N}\Omega^1_{\Sigma/X,p_i}.
\]
On the summand at~$p_i$, the function $1/x$ differs from the scalar
$1/x(\pi(p_i))$ by an element of the maximal ideal
$\mathfrak m_{p_i}$.  The action of $\mathfrak m_{p_i}$ on the finite-length
module $\Omega^1_{\Sigma/X,p_i}$ is nilpotent, so the induced endomorphism has
the form
\[
\frac{1}{x(\pi(p_i))}\mathrm{id}+N_i
\qquad\text{with }N_i\text{ nilpotent}.
\]
Hence the characteristic factor on that summand is
\[
\left(1-\frac{z}{x(\pi(p_i))}\right)^{\ell_i}.
\]
Multiplying over all support points gives the claimed determinant.
\end{proof}

\begin{remark}[Tautological specialization]
\label{rem:tautological-specialization}
Taking the tautological presentation $(F(P),P,t)$ recovers the spectral sheet
field $F(P)$, the spectral cover generated by the visible series itself, and the
branch package of the earlier quadratic calculations.  The present chapter is
therefore not a replacement of the tautological package but its correct
functorial enlargement.
\end{remark}

\begin{definition}[Transport-compatible lifted presentation]
\label{def:transport-compatible-presentation}
Let $\mathcal C$ be a class of spectral objects equipped with visible series
$P_{\cA}(x)$.  A lifted presentation
\[
\cA\longmapsto (L_{\cA},y_{\cA},r_{\cA})
\]
is \emph{transport-compatible} for an exact functor $\Phi\colon\mathcal C\to
\mathcal C$ if for every $\cA$ there exists a rational function
$u_{\Phi,\cA}(x,t)\in F(t)$ such that
\[
y_{\Phi(\cA)}=u_{\Phi,\cA}\bigl(x,y_{\cA}\bigr)
\in L_{\cA}.
\]
\end{definition}

\begin{theorem}[Transport of lifted covers; \ClaimStatusProvedHere]
\label{thm:transport-lifted-covers}
Assume a lifted presentation is transport-compatible for an exact functor
$\Phi$.  Then for each $\cA$ there is a canonical field inclusion
\[
L_{\Phi(\cA)}\hookrightarrow L_{\cA},
\]
hence a canonical finite dominant morphism of normalized covers over~$X$,
\[
f_{\Phi,\cA}\colon \Sigma_{\cA}\to \Sigma_{\Phi(\cA)}.
\]
Moreover:
\begin{enumerate}[label=\textup{(\roman*)}]
\item there is a transitivity triangle
\[
Lf_{\Phi,\cA}^*L_{\Sigma_{\Phi(\cA)}/X}
\longrightarrow
L_{\Sigma_{\cA}/X}
\longrightarrow
L_{\Sigma_{\cA}/\Sigma_{\Phi(\cA)}}
\longrightarrow;
\]
\item the relative cotangent complex
\[
\mathbb D_{\Phi}(\cA):=L_{\Sigma_{\cA}/\Sigma_{\Phi(\cA)}}
\]
measures the exact failure of branch-package preservation;
\item if there is reverse rational transport, or merely
$[L_{\cA}:F]=[L_{\Phi(\cA)}:F]$, then $f_{\Phi,\cA}$ is an isomorphism and the
full lifted branch package is preserved.
\end{enumerate}
\end{theorem}

\begin{proof}
The rational transport law shows that
$y_{\Phi(\cA)}\in F(y_{\cA})=L_{\cA}$, hence
$L_{\Phi(\cA)}=F(y_{\Phi(\cA)})\subseteq L_{\cA}$.  Normalization is
contravariantly functorial for finite field extensions, giving the morphism of
covers.  The triangle in~(i) is the standard transitivity triangle for
cotangent complexes.  If degrees agree, the field inclusion is an equality; if
reverse rational transport exists, equality follows similarly.  A finite
birational morphism between normal curves is an isomorphism.
\end{proof}

\begin{remark}[Defect versus visible invariance]
\label{rem:defect-vs-visible-invariance}
Visible discriminant invariance, as in
Proposition~\ref{prop:ds-package-functoriality}, is weaker than preservation of
the lifted cover.  The determinant formula records only the image of the branch
divisor on~$X$.  The relative cotangent complex $\mathbb D_{\Phi}(\cA)$ records
what visible discriminants cannot: sheet loss, hidden ramification, and the
failure of a lifted coordinate to descend unchanged.
\end{remark}

\section{Quadratic sheet calculus}
\label{sec:quadratic-sheet-calculus}

We first treat the quadratic regime in its sharpest form.  This absorbs the
spectral-sheet-field calculations of the previous notes and identifies the exact
mechanism behind tensor cancellation.

\begin{setup}[Quadratic sheet]
\label{setup:quadratic-sheet}
Let $L=F(u)$ with
\[
u^2=\Delta(x)\in F^{\times}
\]
and $\Delta$ not a square in~$F$.  Every element of~$L$ has a unique
decomposition
\[
y=a+b u,\qquad a,b\in F.
\]
We call $a$ the \emph{even part} and $bu$ the \emph{odd part} of~$y$.
\end{setup}

\begin{definition}[Quadratic spectral datum]
\label{def:quadratic-spectral-datum}
A \emph{quadratic spectral datum} on the sheet~$L/F$ is an element
$y=a+bu\in L$.  It is \emph{purely odd} if $a=0$ and \emph{rational} if
$b=0$.
\end{definition}

\begin{theorem}[Common-sheet multiplication law; \ClaimStatusProvedHere]
\label{thm:quadratic-common-sheet-law}
Let $y_1=a_1+b_1u$ and $y_2=a_2+b_2u$ be quadratic spectral data on the same
sheet~$L/F$.  Then
\[
y_1y_2=(a_1a_2+\Delta b_1b_2)+(a_1b_2+a_2b_1)u.
\]
In particular:
\begin{enumerate}[label=\textup{(\roman*)}]
\item $y_1y_2\in F$ if and only if
\[
a_1b_2+a_2b_1=0;
\]
\item if both factors are purely odd, then
\[
y_1y_2=\Delta b_1b_2\in F.
\]
\end{enumerate}
\end{theorem}

\begin{proof}
This is the multiplication law in the quadratic extension $F(u)$, using
$u^2=\Delta$.  The rationality criterion is exactly the vanishing of the odd
part.
\end{proof}

\begin{theorem}[Local branch survival criterion; \ClaimStatusProvedHere]
\label{thm:local-branch-survival-criterion}
Assume $\Delta\in k[x]$ is squarefree and let $r$ be a simple zero of
$\Delta$.  Suppose $a_i,b_i\in F$ are regular at~$r$.  Then the product
$y_1y_2$ is branched at~$r$ if and only if
\[
a_1(r)b_2(r)+a_2(r)b_1(r)\neq 0.
\]
\end{theorem}

\begin{proof}
Locally one may write
$\Delta(x)=(x-r)u_0(x)$ with $u_0(r)\neq 0$, hence
\[
u=\sqrt{u_0(r)}\,\sqrt{x-r}+O\bigl((x-r)^{3/2}\bigr).
\]
Multiplying the local expansions of $y_1$ and $y_2$ shows that the coefficient
of $\sqrt{x-r}$ in $y_1y_2$ is exactly
$a_1(r)b_2(r)+a_2(r)b_1(r)$.
\end{proof}

\begin{corollary}[Tautological quadratic tensor cancellation; \ClaimStatusProvedHere]
\label{cor:tautological-quadratic-tensor-cancellation}
Assume $P_{\cA}(x)$ and $P_{\cB}(x)$ are tautological quadratic series on the
same spectral sheet and that the visible tensor law is multiplicative,
\[
P_{\cA\otimes\cB}(x)=P_{\cA}(x)P_{\cB}(x).
\]
Then the tensor product is spectrally rational if and only if the odd parts of
$P_{\cA}$ and $P_{\cB}$ satisfy the cancellation relation of
Theorem~\ref{thm:quadratic-common-sheet-law}.  In particular, purely odd
common-sheet tensors become rational.
\end{corollary}

\begin{proof}
Apply Theorem~\ref{thm:quadratic-common-sheet-law} to the tautological lifted
coordinates $y_i=P_{\cA},P_{\cB}$.
\end{proof}

\begin{remark}[What the visible discriminant misses]
\label{rem:visible-discriminant-misses-odd-part}
The common branch polynomial $\Delta$ records where the sheet may branch but not
how the visible series sits on that sheet.  The tensor-cancellation criterion
shows that the missing datum is exactly the odd component of the lifted
coordinate.  This is the first place where the lifted cover, rather than the
visible discriminant, is mathematically unavoidable.
\end{remark}

\section{Distinct quadratic sheets and exact tensor-compositum defects}
\label{sec:distinct-quadratic-sheets}

We next treat the genuinely two-sheet situation.  The resulting classification
is complete: the tensor-compositum field is either quartic or quadratic, and
its ramified defect is governed exactly by intersection of branch divisors.

\begin{setup}[Distinct quadratic sheets]
\label{setup:distinct-quadratic-sheets}
Let
\[
L_A=F(u),\qquad u^2=\Delta_A,
\qquad
L_B=F(v),\qquad v^2=\Delta_B,
\]
with $\Delta_A,\Delta_B\in F^{\times}$ representing distinct square classes.
Set
\[
L:=F(u,v),\qquad w:=uv,
\qquad w^2=\Delta_A\Delta_B.
\]
Then $L/F$ is biquadratic with Galois group
\[
G=\{1,\sigma_u,\sigma_v,\sigma_w=\sigma_u\sigma_v\}\cong (\mathbf Z/2)^2,
\]
where
\[
\sigma_u(u)=-u,\ \sigma_u(v)=v,
\qquad
\sigma_v(u)=u,\ \sigma_v(v)=-v.
\]
\end{setup}

\begin{lemma}[Intermediate fields in the biquadratic case; \ClaimStatusProvedHere]
\label{lem:biquadratic-intermediates}
The proper intermediate fields of $L/F$ are exactly
\[
F(u),\qquad F(v),\qquad F(w).
\]
\end{lemma}

\begin{proof}
By Galois theory, proper intermediate fields correspond to subgroups of index
$2$ in $G$.  The three such subgroups are
$\langle\sigma_u\rangle$, $\langle\sigma_v\rangle$, and
$\langle\sigma_w\rangle$, whose fixed fields are $F(v)$, $F(u)$, and $F(w)$,
respectively.
\end{proof}

\begin{theorem}[Field generated by a biquadratic element; \ClaimStatusProvedHere]
\label{thm:biquadratic-generated-field-classification}
Let
\[
P=\alpha+\beta u+\gamma v+\delta w\in L,
\qquad \alpha,\beta,\gamma,\delta\in F.
\]
Then:
\begin{enumerate}[label=\textup{(\roman*)}]
\item $F(P)=F$ if and only if $\beta=\gamma=\delta=0$;
\item $F(P)=F(u)$ if and only if $\gamma=\delta=0$ and $\beta\neq 0$;
\item $F(P)=F(v)$ if and only if $\beta=\delta=0$ and $\gamma\neq 0$;
\item $F(P)=F(w)$ if and only if $\beta=\gamma=0$ and $\delta\neq 0$;
\item in every other case, $F(P)=L$.
\end{enumerate}
\end{theorem}

\begin{proof}
The field $F(P)$ is the fixed field of the subgroup of $G$ fixing~$P$.
Now
\[
\sigma_u(P)=\alpha-\beta u+\gamma v-\delta w,
\]
so $\sigma_u(P)=P$ if and only if $\beta=\delta=0$.
Similarly,
\[
\sigma_v(P)=P \iff \gamma=\delta=0,
\qquad
\sigma_w(P)=P \iff \beta=\gamma=0.
\]
If all three nonconstant coefficients vanish, then $P\in F$.  If exactly one of
$\beta,\gamma,\delta$ is nonzero, then $P$ is fixed by exactly one nontrivial
involution, so Lemma~\ref{lem:biquadratic-intermediates} identifies the
resulting quadratic field.  In every remaining case no nontrivial involution
fixes $P$, so its stabilizer is trivial and $F(P)=L$.
\end{proof}

\begin{theorem}[Exact quadratic tensor-compositum dichotomy; \ClaimStatusProvedHere]
\label{thm:exact-quadratic-tensor-dichotomy}
Let
\[
y_A=a+bu\in F(u),\qquad y_B=c+dv\in F(v)
\]
with $b,d\neq 0$.
Set $y_{A\otimes B}:=y_Ay_B\in L$.
Then
\[
y_{A\otimes B}=ac+bc\,u+ad\,v+bd\,w.
\]
Consequently,
\[
F\bigl(y_{A\otimes B}\bigr)=
\begin{cases}
F(w),& a=c=0,\\[4pt]
L,& \text{otherwise.}
\end{cases}
\]
Equivalently, in the distinct-sheet quadratic regime:
\begin{enumerate}[label=\textup{(\roman*)}]
\item pure odd $\times$ pure odd lands on the mixed quadratic sheet $F(w)$;
\item every other genuinely quadratic pair generates the full biquadratic
compositum~$L$.
\end{enumerate}
\end{theorem}

\begin{proof}
The displayed expression follows by direct multiplication.  Since $b,d\neq 0$,
the coefficients of $u,v,w$ are respectively $bc$, $ad$, and~$bd$.  If
$a=c=0$, then only the $w$-coefficient survives and
Theorem~\ref{thm:biquadratic-generated-field-classification} gives
$F(y_{A\otimes B})=F(w)$.  Otherwise at least two of the coefficients among
$bc,ad,bd$ are nonzero, so no nontrivial involution fixes the product; hence the
generated field is the full biquadratic compositum~$L$.
\end{proof}

\begin{computation}[Quartic spectral equation in the generic case; \ClaimStatusProvedHere]
\label{comp:quartic-spectral-equation}
Assume $a$ and $c$ are not both zero, so
$F(y_{A\otimes B})=L$ by
Theorem~\ref{thm:exact-quadratic-tensor-dichotomy}.
Write
\[
\alpha:=ac,
\qquad
\beta:=bc,
\qquad
\gamma:=ad,
\qquad
\delta:=bd.
\]
Then $T:=y_{A\otimes B}$ satisfies the depressed quartic equation
\[
(T-\alpha)^4-2S(T-\alpha)^2-8C(T-\alpha)+D=0,
\]
where
\[
S=\beta^2\Delta_A+\gamma^2\Delta_B+\delta^2\Delta_A\Delta_B,
\qquad
C=\beta\gamma\delta\,\Delta_A\Delta_B,
\]
and
\begin{align*}
D={}&
\beta^4\Delta_A^2
+\gamma^4\Delta_B^2
+\delta^4\Delta_A^2\Delta_B^2
-2\beta^2\gamma^2\Delta_A\Delta_B \\
&
-2\beta^2\delta^2\Delta_A^2\Delta_B
-2\gamma^2\delta^2\Delta_A\Delta_B^2.
\end{align*}
\end{computation}

\begin{proof}
The conjugates of $T-\alpha$ under $G$ are
\[
\pm\beta u\pm\gamma v\pm\delta w
\]
with the sign on $w$ equal to the product of the signs on $u$ and~$v$.
Multiplying the four linear factors
\[
\prod_{\varepsilon_1,\varepsilon_2\in\{\pm 1\}}
\Bigl(Z-(\varepsilon_1\beta u+\varepsilon_2\gamma v+
\varepsilon_1\varepsilon_2\delta w)\Bigr)
\]
and substituting $u^2=\Delta_A$, $v^2=\Delta_B$, $w^2=\Delta_A\Delta_B$ gives
the displayed quartic polynomial in $Z=T-\alpha$.
\end{proof}

\begin{definition}[Quadratic tensor defect package]
\label{def:quadratic-tensor-defect-package}
In the setting of Theorem~\ref{thm:exact-quadratic-tensor-dichotomy}, let
\[
K_{A\otimes B}:=F(y_{A\otimes B})\subseteq L.
\]
The \emph{sheet-loss defect} is
\[
\delta_{\otimes}(A,B):=[L:K_{A\otimes B}].
\]
When $K_{A\otimes B}\neq L$, let
\[
q_{A,B}\colon \Sigma_L\to \Sigma_{A\otimes B}
\]
be the induced finite morphism of normalized covers.  The
\emph{ramified tensor defect complex} is
\[
\mathbb D_{\otimes}(A,B):=L_{\Sigma_L/\Sigma_{A\otimes B}}.
\]
\end{definition}

\begin{remark}[Two-layer necessity]
\label{rem:two-layer-necessity}
The number $\delta_{\otimes}(A,B)$ measures total sheet loss.  The cotangent
complex $\mathbb D_{\otimes}(A,B)$ measures only the ramified part of that
loss.  These are genuinely different: in the pure odd $\times$ pure odd case,
$\delta_{\otimes}=2$ while ramification occurs only over common odd branch
points.
\end{remark}

\begin{theorem}[Ramification locus of the pure-odd defect; \ClaimStatusProvedHere]
\label{thm:pure-odd-defect-ramification}
Assume $y_A=bu$ and $y_B=dv$ are pure odd.  Then
$K_{A\otimes B}=F(w)$ and the defect map is
\[
q=q_{A,B}\colon \Sigma_L\to \Sigma_w.
\]
Let $x_0\in X$ be a closed point and define parity data
\[
m(x_0):=v_{x_0}(\Delta_A)\bmod 2,
\qquad
n(x_0):=v_{x_0}(\Delta_B)\bmod 2.
\]
Then $q$ ramifies above $x_0$ if and only if
\[
m(x_0)=n(x_0)=1.
\]
Equivalently, the ramified tensor defect is supported exactly on the
intersection of the odd branch divisors of $\Delta_A$ and~$\Delta_B$.
\end{theorem}

\begin{proof}
The field statement $K_{A\otimes B}=F(w)$ is
Theorem~\ref{thm:exact-quadratic-tensor-dichotomy}.  To test ramification,
complete at~$x_0$.  Because $k$ is algebraically closed, every unit in the
completed local ring has a square root; hence only the parity of the valuation
matters.

If $(m,n)=(0,0)$, then both quadratic extensions are unramified and so is~$q$.
If $(m,n)=(1,0)$, then locally $u$ is ramified while $v$ is unramified, and
$w=uv$ has the same odd parity as~$u$; therefore $F(w)$ already carries the
ramification and $L/F(w)$ is unramified.  The case $(0,1)$ is identical.
Finally, if $(m,n)=(1,1)$ then $u$ and $v$ are both locally ramified but
$w^2=\Delta_A\Delta_B$ has even valuation, so $F(w)/F$ is unramified while
$L/F(w)$ remains quadratic ramified.  Hence $q$ ramifies precisely in that
case.
\end{proof}

\begin{corollary}[Determinant of the pure-odd defect operator; \ClaimStatusProvedHere]
\label{cor:pure-odd-defect-operator}
Assume in addition that $\Delta_A$ and $\Delta_B$ are squarefree and unramified
over $x=0$.  Let
\[
B_A^{\mathrm{fin}},\ B_B^{\mathrm{fin}}\subset X\setminus\{0,\infty\}
\]
be the finite odd branch sets of $\Delta_A$ and $\Delta_B$.
Then multiplication by $1/x$ defines an operator
\[
T^{\mathrm{def}}_{A,B}
\in \operatorname{End}_k H^0\bigl(\Sigma_L,\Omega^1_{\Sigma_L/\Sigma_w}\bigr)
\]
with characteristic polynomial
\[
\det\bigl(1-zT^{\mathrm{def}}_{A,B}\bigr)
=
\prod_{r\in B_A^{\mathrm{fin}}\cap B_B^{\mathrm{fin}}}
\left(1-\frac{z}{r}\right)^2.
\]
Consequently,
\[
\dim_k H^0\bigl(\Sigma_L,\Omega^1_{\Sigma_L/\Sigma_w}\bigr)
=2\#\bigl(B_A^{\mathrm{fin}}\cap B_B^{\mathrm{fin}}\bigr).
\]
\end{corollary}

\begin{proof}
By Theorem~\ref{thm:pure-odd-defect-ramification}, a finite point
$r\in X$ contributes only when it lies in the common odd branch set.  At such a
point $\Sigma_w\to X$ is unramified, hence there are two points of $\Sigma_w$
above~$r$; over each of them the quadratic map $q$ is tamely ramified of index
$2$, so the local model is $k[[t]]\subset k[[u]]$ with $t=u^2$.  Therefore each
ramified point contributes a length-$1$ summand to
$\Omega^1_{\Sigma_L/\Sigma_w}$, and there are two such points above each common
branch value~$r$.  The determinant formula is then
Theorem~\ref{thm:determinantal-branch-formula-general} applied to $q$.
\end{proof}

\begin{theorem}[Complete quadratic tensor calculus; \ClaimStatusProvedHere]
\label{thm:complete-quadratic-tensor-calculus}
For two genuinely quadratic spectral data, the degree of the tensor-compositum
field is completely determined as follows.
\begin{enumerate}[label=\textup{(\roman*)}]
\item \emph{Same quadratic sheet.}
If both data lie on the same sheet $F(u)$, then the tensor product has degree
$1$ over~$F$ exactly on the cancellation locus of
Theorem~\ref{thm:quadratic-common-sheet-law}; otherwise it has degree~$2$.
\item \emph{Distinct quadratic sheets.}
If the two sheets are distinct, then pure odd $\times$ pure odd has degree~$2$,
while every other genuinely quadratic pair has degree~$4$.
\end{enumerate}
Thus the quadratic frontier is completely classified.
\end{theorem}

\begin{proof}
Part~(i) is Theorem~\ref{thm:quadratic-common-sheet-law}; part~(ii) is
Theorem~\ref{thm:exact-quadratic-tensor-dichotomy}.
\end{proof}

\section{Cubic spectral packets and the resolvent tower}
\label{sec:cubic-resolvent-tower}

The correct cubic analogue of the quadratic sheet is no longer a single cover
but a two-stage tower: a quadratic odd cover recording transposition branching,
and a cyclic cubic even cover over it recording $3$-cycle branching.

\begin{setup}[Depressed cubic spectral datum]
\label{setup:depressed-cubic}
Assume $\mu_3\subset k$.  Let
\[
f(y)=y^3+p(x)y+q(x)\in F[y]
\]
be irreducible and separable, and set
\[
\Delta_f:=-4p^3-27q^2\in F^{\times}.
\]
Assume $\Delta_f\notin F^{\times 2}$, so the Galois closure has group~$S_3$.
Let $K/F$ be the cubic field generated by a root of~$f$.
Define the quadratic resolvent field
\[
Q:=F(s),\qquad s^2=q^2+\frac{4}{27}p^3=-\frac{1}{27}\Delta_f,
\]
and the Cardano parameter
\[
\theta:=\frac{-q+s}{2}\in Q^{\times}.
\]
Finally set
\[
L:=Q(u),\qquad u^3=\theta.
\]
\end{setup}

\begin{proposition}[Cardano identities; \ClaimStatusProvedHere]
\label{prop:cardano-identities}
In the setting of Setup~\ref{setup:depressed-cubic} one has
\[
\iota(\theta)=\frac{-q-s}{2}= -\frac{p^3}{27\theta},
\]
where $\iota\in\operatorname{Gal}(Q/F)$ is the nontrivial involution.  In particular,
\[
\iota([\theta])=[\theta]^{-1}
\qquad\text{in }Q^{\times}/Q^{\times 3}.
\]
\end{proposition}

\begin{proof}
The product of the two Cardano factors is
\[
\frac{-q+s}{2}\cdot\frac{-q-s}{2}=\frac{q^2-s^2}{4}
=-\frac{p^3}{27}.
\]
Hence $\iota(\theta)=(-q-s)/2=-p^3/(27\theta)$.  Since $-p^3/27$ is a cube in
$Q^{\times}$, the class relation follows.
\end{proof}

\begin{theorem}[Cardano tower and cubic quotients; \ClaimStatusProvedHere]
\label{thm:cardano-tower-cubic-quotients}
In the setting of Setup~\ref{setup:depressed-cubic}, let $\zeta_3\in k$ be a
primitive cube root of unity and define
\[
y_i:=\zeta_3^i u-\frac{p}{3\zeta_3^i u},
\qquad i\in \mathbf Z/3.
\]
Then:
\begin{enumerate}[label=\textup{(\roman*)}]
\item each $y_i$ is a root of~$f$;
\item $L/F$ is the Galois closure of~$K/F$ and
\[
\operatorname{Gal}(L/F)\cong S_3,
\qquad
\operatorname{Gal}(L/Q)\cong C_3;
\]
\item the three cubic subfields of $L/F$ are
\[
K_i:=F(y_i),\qquad i\in \mathbf Z/3,
\]
and the set
\[
\mathfrak T_f:=\{K_0,K_1,K_2\}
\]
is a canonical $C_3$-torsor.
\end{enumerate}
\end{theorem}

\begin{proof}
Set $v_i:=-p/(3\zeta_3^i u)$.  Then
$y_i=\zeta_3^i u+v_i$ and
\[
(\zeta_3^i u)^3+v_i^3
=\theta-\frac{p^3}{27\theta}
=-q,
\qquad
(\zeta_3^i u)\,v_i=-\frac{p}{3}.
\]
Therefore
\[
y_i^3+py_i+q
=(\zeta_3^i u+v_i)^3+p(\zeta_3^i u+v_i)+q=0,
\]
proving~(i).  Since $\Delta_f$ is not a square, the Galois group of the
irreducible cubic is not contained in~$A_3$, hence is~$S_3$; the subgroup over
$Q$ is the order-$3$ subgroup generated by $u\mapsto \zeta_3 u$.  The three
cubic subfields correspond to the three order-$2$ subgroups of~$S_3$, and the
normal subgroup $A_3\cong C_3$ acts simply transitively on them by conjugation,
which gives the torsor structure.
\end{proof}

\begin{definition}[Cubic spectral packet]
\label{def:cubic-spectral-packet}
The \emph{unpointed cubic spectral packet} attached to~$f$ is the triple
\[
\mathfrak S^{S_3}(f):=
\bigl(Q/F,\ \ell_f,\ \mathfrak T_f\bigr),
\]
where
\[
\ell_f:=\langle[\theta]\rangle\subset Q^{\times}/Q^{\times 3}
\]
is the $1$-dimensional Kummer line generated by the Cardano class, and
$\mathfrak T_f$ is the $C_3$-torsor of cubic quotients from
Theorem~\ref{thm:cardano-tower-cubic-quotients}.  The \emph{pointed cubic
packet} is the refinement
\[
\mathfrak S^{S_3}_{\mathrm{pt}}(f):=
\bigl(Q/F,\ [\theta],\ \mathfrak T_f\bigr).
\]
\end{definition}

\begin{remark}[Why the line is unpointed and the class is pointed]
\label{rem:line-vs-class}
The cyclic cubic extension $L/Q$ depends only on the line $\ell_f$ because
$Q(\sqrt[3]{a})=Q(\sqrt[3]{a^{-1}})$ and both classes generate the same order-$3$
subgroup of $Q^{\times}/Q^{\times 3}$.  The pointed class $[\theta]$ remembers
which Cardano orientation has been chosen on the quadratic resolvent sheet.
This distinction is exactly what was missing from the earlier, overcompressed
formulation of the cubic defect.
\end{remark}

\begin{construction}[Odd and even covers]
\label{cons:odd-even-covers}
Let
\[
\Sigma^{\mathrm{odd}}_f\to X
\]
be the normalization of $X$ in the quadratic resolvent field~$Q$, and let
\[
\widetilde\Sigma_f\to \Sigma^{\mathrm{odd}}_f
\]
be the normalization in the cyclic cubic field~$L=Q(u)$.  Define the
corresponding branch complexes
\[
\mathbb B^{\mathrm{odd}}_f:=L_{\Sigma^{\mathrm{odd}}_f/X},
\qquad
\mathbb B^{\mathrm{even}}_f:=L_{\widetilde\Sigma_f/\Sigma^{\mathrm{odd}}_f}.
\]
If both stages are unramified over $x=0$, define branch operators by
multiplication by $1/x$ on the global sections of the corresponding sheaves of
relative differentials.
\end{construction}

\begin{theorem}[Odd-even branch splitting; \ClaimStatusProvedHere]
\label{thm:odd-even-branch-splitting}
In the setting of Construction~\ref{cons:odd-even-covers}, there is a
transitivity triangle
\[
q^*L_{\Sigma^{\mathrm{odd}}_f/X}
\longrightarrow
L_{\widetilde\Sigma_f/X}
\longrightarrow
L_{\widetilde\Sigma_f/\Sigma^{\mathrm{odd}}_f}
\longrightarrow,
\]
where $q\colon\widetilde\Sigma_f\to\Sigma^{\mathrm{odd}}_f$ is the cyclic cubic
cover.  Moreover, the ramification support splits by inertia type:
\begin{enumerate}[label=\textup{(\roman*)}]
\item points with transposition inertia contribute only to
$\mathbb B^{\mathrm{odd}}_f$;
\item points with $3$-cycle inertia contribute only to
$\mathbb B^{\mathrm{even}}_f$.
\end{enumerate}
\end{theorem}

\begin{proof}
The triangle is the standard transitivity triangle.  For the inertia analysis,
let $p\in\widetilde\Sigma_f$ with inertia subgroup $I_p\subset S_3$.  Because we
work over characteristic~$0$, inertia is cyclic.  The only nontrivial cyclic
subgroups of $S_3$ are transpositions and $A_3$.  If $I_p$ is a transposition,
its image in $S_3/A_3\cong C_2$ is nontrivial, so the quadratic resolvent cover
ramifies, while $I_p\cap A_3=1$ so the cubic stage is étale.  If
$I_p=A_3$, then the image in $S_3/A_3$ is trivial, so the odd cover is étale,
whereas the cubic stage is totally ramified of degree~$3$.
\end{proof}

\begin{corollary}[Local lengths in the tame simple case; \ClaimStatusProvedHere]
\label{cor:odd-even-local-lengths}
Assume all branch points are tame and simple.  Then:
\begin{enumerate}[label=\textup{(\roman*)}]
\item at a transposition branch point the odd branch sheaf has length~$1$;
\item at a $3$-cycle branch point the even branch sheaf has length~$2$.
\end{enumerate}
Consequently,
\begin{align*}
\det\bigl(1-zT_f^{\mathrm{odd}}\bigr)
&=
\prod_{q\in \operatorname{Supp}\Omega^1_{\Sigma^{\mathrm{odd}}_f/X}}
\left(1-\frac{z}{x(q)}\right)^{\ell_q},\\
\det\bigl(1-zT_f^{\mathrm{even}}\bigr)
&=
\prod_{p\in \operatorname{Supp}\Omega^1_{\widetilde\Sigma_f/\Sigma^{\mathrm{odd}}_f}}
\left(1-\frac{z}{x(p)}\right)^{m_p},
\end{align*}
with $\ell_q=1$ and $m_p=2$ on simple branch points.
\end{corollary}

\begin{proof}
For a simple quadratic branch point the local model is
$k[[t]]\subset k[[u]]$, $t=u^2$, so
\[
\Omega^1_{k[[u]]/k[[t]]}\cong k[[u]]/(u)\,du,
\]
which has length~$1$.  For a simple cubic Kummer branch point the local model
is $k[[s]]\subset k[[v]]$, $s=v^3$, so
\[
\Omega^1_{k[[v]]/k[[s]]}\cong k[[v]]/(v^2)\,dv,
\]
which has length~$2$.  The determinant formulas are
Theorem~\ref{thm:determinantal-branch-formula-general} applied to the odd and
even stages separately.
\end{proof}

\section{Comparison of cubic packets: Kummer lines, pointed classes, and defect covers}
\label{sec:comparison-cubic-packets}

We now classify exactly what it means for two cubic packets to agree.  The
answer has three layers: the quadratic odd sheet, the unpointed Kummer line, and
the residual $C_3$-torsor of cubic quotients.

\begin{setup}[Two cubic packets with common odd sheet]
\label{setup:two-cubic-packets}
Let $f_A$ and $f_B$ be two depressed irreducible cubic equations over~$F$ with
quadratic resolvent fields identified as
\[
Q_A\cong Q_B=:Q.
\]
Write $\theta_A,\theta_B\in Q^{\times}$ for the corresponding pointed Cardano
parameters and
\[
\ell_A:=\langle[\theta_A]\rangle,
\qquad
\ell_B:=\langle[\theta_B]\rangle
\subset Q^{\times}/Q^{\times 3}
\]
for the corresponding Kummer lines.  Let
\[
L_A:=Q(\sqrt[3]{\theta_A}),
\qquad
L_B:=Q(\sqrt[3]{\theta_B})
\]
be the cyclic cubic closures over~$Q$.
\end{setup}

\begin{theorem}[Unpointed equality criterion; \ClaimStatusProvedHere]
\label{thm:unpointed-kummer-equality}
In Setup~\ref{setup:two-cubic-packets}, the following are equivalent:
\begin{enumerate}[label=\textup{(\roman*)}]
\item $L_A\cong_Q L_B$ as cyclic cubic extensions of~$Q$;
\item $\ell_A=\ell_B$ as lines in $Q^{\times}/Q^{\times 3}$;
\item $[\theta_B]$ equals either $[\theta_A]$ or $[\theta_A]^{-1}$ in
$Q^{\times}/Q^{\times 3}$.
\end{enumerate}
\end{theorem}

\begin{proof}
Kummer theory identifies cyclic cubic extensions of~$Q$ with $1$-dimensional
subspaces of the $\mathbf F_3$-vector space $Q^{\times}/Q^{\times 3}$.  This
proves the equivalence of~(i) and~(ii).  Since a nonzero $1$-dimensional
$\mathbf F_3$-subspace has exactly two nonzero elements, namely a class and its
inverse,~(ii) and~(iii) are equivalent.
\end{proof}

\begin{theorem}[Pointed equality criterion; \ClaimStatusProvedHere]
\label{thm:pointed-kummer-equality}
Keep the pointed Cardano classes from Setup~\ref{setup:two-cubic-packets}.  The
pointed cyclic cubic extensions are equal if and only if
\[
[\theta_B/\theta_A]=1
\qquad\text{in }Q^{\times}/Q^{\times 3}.
\]
Equivalently, the pointed Cardano classes coincide.
\end{theorem}

\begin{proof}
This is the standard Kummer criterion for equality of pointed cyclic cubic
extensions: $Q(\sqrt[3]{a})=Q(\sqrt[3]{b})$ with chosen generators agreeing if
and only if $b/a$ is a cube in~$Q^{\times}$.
\end{proof}

\begin{remark}[Why both criteria are needed]
\label{rem:why-both-criteria-are-needed}
The unpointed criterion classifies the cyclic cubic extension over the odd
sheet.  The pointed criterion is stronger: it remembers which Cardano generator
on that cyclic cubic extension has been chosen.  Earlier formulations collapsed
these two levels and thereby misstated the true equality criterion.  The present
package separates them cleanly.
\end{remark}

\begin{proposition}[Selector torsor defect; \ClaimStatusProvedHere]
\label{prop:selector-torsor-defect}
Assume $L_A\cong_Q L_B$ and fix a $Q$-isomorphism
\[
\phi\colon L_A\xrightarrow{\sim} L_B.
\]
Then there exists a unique element
\[
\tau(A,B)\in C_3
\]
such that the induced bijection on cubic quotients is translation by
$\tau(A,B)$:
\[
\phi\bigl(K_{A,i}\bigr)=K_{B,i+\tau(A,B)}
\qquad\text{for all }i\in\mathbf Z/3.
\]
Thus equality of cyclic closures leaves a residual $C_3$-torsorial defect among
the cubic quotients.
\end{proposition}

\begin{proof}
By Theorem~\ref{thm:cardano-tower-cubic-quotients}, each packet carries a
canonical $C_3$-torsor of cubic subfields.  Any $Q$-isomorphism of the cyclic
closures induces a bijection between the two torsors, and any bijection between
principal homogeneous $C_3$-spaces is translation by a unique element of~$C_3$.
\end{proof}

\begin{proposition}[Base change to the odd cover; \ClaimStatusProvedHere]
\label{prop:base-change-to-odd-cover}
For a fixed cubic packet~$f$, every cubic quotient becomes the same cover after
base change to the odd resolvent cover:
\[
\operatorname{Norm}\Bigl(\Sigma_{K_i}\times_X\Sigma^{\mathrm{odd}}_f\Bigr)
\cong
\widetilde\Sigma_f.
\]
\end{proposition}

\begin{proof}
One has $[K_i:F]=3$ and $[Q:F]=2$, so $K_i\cap Q=F$.  Hence the compositum
$K_iQ$ has degree $6$ over~$F$, which equals $[L:F]$; therefore $K_iQ=L$.
Normalization of the fiber product corresponds to taking the compositum of
function fields.
\end{proof}

\begin{definition}[Pointed cubic defect cover]
\label{def:pointed-cubic-defect-cover}
Assume the odd sheet~$Q$ has been identified and the pointed Cardano classes
$\theta_A,\theta_B$ have been fixed.  Set
\[
\xi:=\theta_B/\theta_A\in Q^{\times}
\]
and define the \emph{pointed relative cubic defect field}
\[
M_{A,B}:=Q(w),\qquad w^3=\xi.
\]
Let
\[
\rho_{A,B}\colon \Sigma_{\xi}\to \Sigma_Q
\]
be the corresponding normalized cover.
\end{definition}

\begin{remark}[Dependence only on the unpointed relative line]
\label{rem:defect-cover-inversion-invariance}
Replacing $\xi$ by $\xi^{-1}$ does not change the extension $Q(w)/Q$, since
$Q(\sqrt[3]{\xi^{-1}})=Q(1/w)=Q(w)$.  Hence the defect cover depends only on the
line generated by the pointed ratio~$[\xi]$.
\end{remark}

\begin{theorem}[Local cubic ramification criterion; \ClaimStatusProvedHere]
\label{thm:local-cubic-ramification-criterion}
Let $p\in \Sigma_Q$ be a closed point and write
\[
\bar v_p(\xi):=v_p(\xi)\bmod 3.
\]
Then the pointed defect cover $\rho_{A,B}$ ramifies above $p$ if and only if
\[
\bar v_p(\xi)\neq 0.
\]
In that case the ramification is tame and totally ramified of degree~$3$, and
for every point $q$ above~$p$ one has
\[
\operatorname{length}_{\mathcal O_{\Sigma_{\xi},q}}
\Omega^1_{\Sigma_{\xi}/\Sigma_Q,q}=2.
\]
\end{theorem}

\begin{proof}
Complete at~$p$ and write
\[
\xi=\pi^{m}u_0,
\]
with $\pi$ a uniformizer and $u_0$ a unit.  Since $k$ is algebraically closed,
$u_0$ has a cube root in the completed local field, so the extension is locally
isomorphic to
\[
K\subset K[t]/(t^3-\pi^m).
\]
If $m\equiv 0\pmod 3$, then $\pi^m$ is a cube and the extension is étale.  If
$m\not\equiv 0\pmod 3$, the extension is tamely totally ramified of degree~$3$.
The local model is then $s=t^3$, so
\[
\Omega^1_{k[[t]]/k[[s]]}\cong k[[t]]/(t^2)\,dt,
\]
which has length~$2$.
\end{proof}

\begin{definition}[Pointed cubic defect complex]
\label{def:pointed-cubic-defect-complex}
The \emph{pointed cubic defect complex} is
\[
\mathbb D^{(3)}_{\mathrm{pt}}(A,B):=L_{\Sigma_{\xi}/\Sigma_Q},
\]
and the corresponding branch space is
\[
V^{(3)}_{\mathrm{pt}}(A,B):=
H^0\bigl(\Sigma_{\xi},\Omega^1_{\Sigma_{\xi}/\Sigma_Q}\bigr).
\]
\end{definition}

\begin{corollary}[Determinant of the pointed cubic defect operator; \ClaimStatusProvedHere]
\label{cor:pointed-cubic-defect-operator}
Assume $\rho_{A,B}$ is unramified above the points of $\Sigma_Q$ lying over
$x=0$.  Then multiplication by $1/x$ defines an operator
\[
T^{(3)}_{\mathrm{pt}}(A,B)
\in \operatorname{End}_k V^{(3)}_{\mathrm{pt}}(A,B)
\]
with characteristic polynomial
\[
\det\bigl(1-zT^{(3)}_{\mathrm{pt}}(A,B)\bigr)
=
\prod_{p\in \operatorname{Ram}(\rho_{A,B})}
\left(1-\frac{z}{x(p)}\right)^2.
\]
If all ramification points are simple and finite, then
\[
\dim_k V^{(3)}_{\mathrm{pt}}(A,B)
=2\#\operatorname{Ram}(\rho_{A,B}).
\]
\end{corollary}

\begin{proof}
Apply Theorem~\ref{thm:determinantal-branch-formula-general} to the finite map
$\rho_{A,B}$ and use Theorem~\ref{thm:local-cubic-ramification-criterion} to
identify each local length with~$2$.
\end{proof}

\begin{corollary}[No étale cubic defect over a rational odd sheet; \ClaimStatusProvedHere]
\label{cor:no-etale-cubic-defect-rational-sheet}
Assume $\Sigma_Q\cong \mathbf P^1$.  Then every nontrivial pointed cubic defect
cover $\rho_{A,B}$ ramifies somewhere.  Equivalently, if
$\rho_{A,B}$ is everywhere étale, then it is trivial.
\end{corollary}

\begin{proof}
Finite étale cyclic cubic covers of $\mathbf P^1$ correspond to $3$-torsion line
bundles on $\mathbf P^1$, but
$\operatorname{Pic}(\mathbf P^1)\cong \mathbf Z$ has no nontrivial $3$-torsion.
\end{proof}

\begin{corollary}[Genus of the pointed cubic defect cover; \ClaimStatusProvedHere]
\label{cor:genus-pointed-cubic-defect}
Assume $\Sigma_Q\cong \mathbf P^1$ and the pointed defect cover is ramified at
exactly $r$ points, all totally ramified of degree~$3$.  Then
\[
g(\Sigma_{\xi})=r-2.
\]
In particular, the first nontrivial pointed cubic defect has $r=2$ and genus
$0$, while the first elliptic defect occurs at $r=3$.
\end{corollary}

\begin{proof}
Riemann--Hurwitz gives
\[
2g(\Sigma_{\xi})-2
=3\bigl(2g(\mathbf P^1)-2\bigr)+r(3-1)
=-6+2r,
\]
so $g(\Sigma_{\xi})=r-2$.
\end{proof}

\section{Rational shadows and the \texorpdfstring{$\widehat{\mathfrak{sl}}_3\rightsquigarrow\mathcal W_3$}{sl3 -> W3} frontier}
\label{sec:rational-shadows-sl3-w3}

The present repository predicts rational visible series for the
$\widehat{\mathfrak{sl}}_3$ and $\mathcal W_3$ orbit
(Conjectures~\ref{conj:sl3-bar-gf} and~\ref{conj:w3-bar-gf}), both sharing the
quadratic visible factor $1-3x-x^2$; see also the discussion in
Appendix~\ref{sec:discriminant-families}.  The tautological spectral cover is
therefore trivial in that visible model.  The cubic package constructed above is
not a claim that those visible rational functions are themselves cubic over~$F$.
It is a claim about a \emph{hidden lift} above their common visible shadow.

\begin{definition}[Hidden cubic lift of a rational-shadow pair]
\label{def:hidden-cubic-lift-rational-pair}
Let $P_1(x),P_2(x)\in F$ be rational visible series sharing a common visible
quadratic factor $D(x)\in F$.  A \emph{hidden cubic lift} of the pair consists
of the following data.
\begin{enumerate}[label=\textup{(\roman*)}]
\item the common odd cover
\[
Q:=F(s),\qquad s^2=D(x);
\]
\item two pointed Cardano classes
$\theta_1,\theta_2\in Q^{\times}$, defining pointed cubic packets over~$Q$;
\item rational reconstruction laws expressing $P_1$ and $P_2$ as rational
functions of chosen cubic coordinates on the corresponding cubic quotients.
\end{enumerate}
\end{definition}

\begin{remark}[Visible invariance versus lifted invariance]
\label{rem:visible-vs-lifted-invariance-sl3}
A hidden cubic lift is compatible with visible rationality because the visible
series are allowed to be rational expressions in the lifted cubic coordinate.
Thus visible equality of the quadratic factor $D(x)$ detects only the odd sheet.
The even Kummer line and the selector torsor are invisible at that level.
\end{remark}

\begin{theorem}[Exact defect package for a hidden cubic lift; \ClaimStatusProvedHere]
\label{thm:cubic-rational-shadow-frontier}
Assume the conjectural rational-shadow pair
$\bigl(P_{\widehat{\mathfrak{sl}}_3},P_{\mathcal W_3}\bigr)$ admits a hidden cubic
lift over the common odd sheet
\[
Q=F(s),\qquad s^2=1-3x-x^2.
\]
Then the full lifted DS defect is determined by the following three pieces of
data.
\begin{enumerate}[label=\textup{(\roman*)}]
\item the common odd cover $Q/F$, which is already visible from the shared factor
$1-3x-x^2$;
\item the unpointed even defect, namely the comparison of the Kummer lines
$\ell_{\widehat{\mathfrak{sl}}_3}$ and $\ell_{\mathcal W_3}$;
\item if those lines agree, the residual selector defect
\[
\tau\bigl(\widehat{\mathfrak{sl}}_3,\mathcal W_3\bigr)\in C_3.
\]
\end{enumerate}
If, in addition, the pointed Cardano classes are chosen on the common odd sheet,
then the pointed defect cover is
\[
Q(w),\qquad
w^3=\theta_{\mathcal W_3}/\theta_{\widehat{\mathfrak{sl}}_3}.
\]
Because the odd curve $s^2=1-3x-x^2$ is rational, any nontrivial pointed defect
must ramify by Corollary~\ref{cor:no-etale-cubic-defect-rational-sheet}; if it
ramifies at $r$ points, its normalization has genus $r-2$ by
Corollary~\ref{cor:genus-pointed-cubic-defect}.
\end{theorem}

\begin{proof}
The hidden cubic lift supplies two pointed cubic packets over the same odd sheet,
so the exact comparison theory of
Section~\ref{sec:comparison-cubic-packets} applies.  The visible factor
$1-3x-x^2$ determines the odd sheet and nothing beyond it.  The unpointed even
comparison is Theorem~\ref{thm:unpointed-kummer-equality}; the residual selector
defect is Proposition~\ref{prop:selector-torsor-defect}.  The pointed relative
cover is Definition~\ref{def:pointed-cubic-defect-cover}.  Since the odd curve
$s^2=1-3x-x^2$ is a smooth double cover of $\mathbf P^1$ branched at two points,
it is rational, so Corollaries
\ref{cor:no-etale-cubic-defect-rational-sheet} and
\ref{cor:genus-pointed-cubic-defect} finish the argument.
\end{proof}

\begin{remark}[What the next computation actually is]
\label{rem:what-next-computation-actually-is}
The next exact computation on the
$\widehat{\mathfrak{sl}}_3\rightsquigarrow\mathcal W_3$ frontier is therefore
not ``find another visible denominator.''  It is:
\begin{enumerate}[label=\textup{(\roman*)}]
\item construct a hidden cubic lift over the rational odd sheet
$s^2=1-3x-x^2$;
\item extract the Kummer line (or pointed Cardano class) of the cyclic cubic
closure over that sheet;
\item compute its defect divisor against the DS image;
\item determine whether any remaining discrepancy is only the selector torsor.
\end{enumerate}
That is the precise cubic $S_3$ analogue of the quadratic tensor-defect package.
\end{remark}

\section{Hierarchy theorem and summary}
\label{sec:hierarchy-summary}

\begin{theorem}[Hierarchy of spectral invariants; \ClaimStatusProvedHere]
\label{thm:hierarchy-of-spectral-invariants}
For a spectral object with a chosen lifted presentation, the invariants discussed
in this chapter form a strict hierarchy:
\begin{enumerate}[label=\textup{(\roman*)}]
\item the \emph{visible shadow}: denominator/discriminant on~$X$;
\item the \emph{lifted cover}: the normalized finite cover $\Sigma\to X$;
\item in the quadratic case, the \emph{odd component of the lifted coordinate},
which controls common-sheet tensor cancellation;
\item in the distinct quadratic case, the \emph{two-layer tensor defect}
$\bigl(\delta_{\otimes},\mathbb D_{\otimes}\bigr)$;
\item in the cubic $S_3$ case, the \emph{odd sheet}, the \emph{even Kummer
line}, and the \emph{selector torsor}.
\end{enumerate}
Each successive layer contains information invisible to the previous ones, and
none of the later tensor- and transport-theoretic formulas can be recovered from
the visible shadow alone.
\end{theorem}

\begin{proof}
The visible shadow is recovered from the branch operator by
Theorem~\ref{thm:determinantal-branch-formula-general}, but
Remark~\ref{rem:visible-discriminant-misses-odd-part} shows that it does not
determine the odd part of a quadratic lifted coordinate.  Theorem
\ref{thm:complete-quadratic-tensor-calculus} shows that this odd part controls
quadratic tensor cancellation.  Definition~\ref{def:quadratic-tensor-defect-package}
and Theorem~\ref{thm:pure-odd-defect-ramification} show that sheet loss and
ramified defect are different invariants in the distinct-sheet case.  Finally,
Definitions~\ref{def:cubic-spectral-packet} and
\ref{def:pointed-cubic-defect-cover}, together with
Theorems~\ref{thm:unpointed-kummer-equality} and
\ref{thm:pointed-kummer-equality}, show that the cubic $S_3$ regime contains two
new layers beyond the odd resolvent cover: the even Kummer line and the selector
torsor.
\end{proof}

\begin{corollary}[Platonic form of the frontier; \ClaimStatusProvedHere]
\label{cor:platonic-form-frontier}
The exact frontier after the current quadratic results is this.
\begin{enumerate}[label=\textup{(\roman*)}]
\item The quadratic tensor calculus is complete.
\item Any genuinely new phenomenon in the next rank must be cubic or higher.
\item In the first cubic case, the problem is no longer to detect the visible
odd sheet --- that is already the discriminant problem --- but to extract and
compare the even Kummer line and the selector torsor above it.
\end{enumerate}
\end{corollary}

\begin{proof}
Part~(i) is Theorem~\ref{thm:complete-quadratic-tensor-calculus}.  Part~(ii)
follows because quadratic phenomena are exhausted by the same-sheet and
distinct-sheet classifications.  Part~(iii) is exactly the content of
Theorem~\ref{thm:cubic-rational-shadow-frontier}.
\end{proof}

\begin{remark}[Integration point in the monograph]
\label{rem:integration-point-monograph}
This chapter is designed to sit between the visible discriminant package of
Theorem~\ref{thm:discriminant-spectral} and the programme layer around the
$\widehat{\mathfrak{sl}}_3/\mathcal W_3$ orbit.  Its S-level statements are
proved in full generality.  Its application to the actual monograph frontier is
conditional only at the point where a hidden cubic lift must still be extracted
from the bar/BRST data.
\end{remark}
